\part{Six-Vertex Model}

\section{Row-to-row transfer matrices}

\begin{exercise}[Ice configurations and the row transfer operator]
Let $S=\{+,-\}$ label arrows pointing up or right, and down or left,
respectively, and put $V=\C^S$ with counting inner product. At each
square-lattice vertex require two incoming and two outgoing arrows.
Regard the left and bottom arrows as inputs and the right and top arrows
as outputs. The vertex operator on $V\otimes V$ has kernel
\[
 R^{++}_{++}=R^{--}_{--}=a,\quad
 R^{+-}_{+-}=R^{-+}_{-+}=b,\quad
 R^{-+}_{+-}=R^{+-}_{-+}=c,
\]
with all other values zero. For $a,b,c>0$ the weight of a configuration
is the product of its vertex weights; normalize these weights by their
sum to obtain its Gibbs measure. On a cylinder of circumference $L$ set
$\mathcal H_L=V^{\otimes L}$ and
\[
 \mathcal T_0=R_{0L}\cdots R_{01},\qquad
 T=\Tr_0\mathcal T_0.
\]
Here $0$ is the horizontal arrow space and $\Tr_0$ is the partial trace
of the section Transfer matrices.
\begin{enumerate}[label=(\alph*)]
\item Expand the partial trace as a sum over horizontal arrow
configurations and prove that $Z_{L,N}=\Tr(T^N)$ on the $L$ by $N$ torus.
\item Let $(Zf)(s)=s f(s)$, with $+=1$, $-=-1$, and
$(Xf)(s)=f(-s)$. Prove that $T$ commutes with
$S^z_{\rm tot}=\frac12\sum_j Z_j$, with $X^{\otimes L}$, and with
cyclic translation. Decompose $\mathcal H_L$ into the sectors
$\mathcal H_{L,M}$ with exactly $M$ down arrows.
\item Prove that the two polarized sectors have transfer eigenvalue
$a^L+b^L$. Show that multiplying all vertex weights by $\rho>0$
adds $\log\rho$ to the pressure per vertex.
\end{enumerate}
\end{exercise}

\section{Spectral parameter}

\begin{exercise}[The trigonometric vertex operator and anisotropy]
For $\eta>0$ and a complex spectral parameter $u$ use the preceding
kernel with
\[
 a(u)=\sinh(u+\eta),\qquad b(u)=\sinh u,\qquad c(u)=\sinh\eta.
\]
The anisotropy is $\Delta=(a^2+b^2-c^2)/(2ab)$ wherever $ab\ne0$.
Let $P(v\otimes w)=w\otimes v$.
\begin{enumerate}[label=(\alph*)]
\item Prove $\Delta=\cosh\eta$, $R(0)=\sinh\eta\,P$, and
\[
 R(u)R(-u)=\bigl(\sinh^2\eta-\sinh^2u\bigr)\id.
\]
Determine the exceptional parameters where $R(u)$ is singular.
\item Continue to $(u,\eta)=(\ii v,\ii\gamma)$ and divide all
weights by $\ii$. Prove that the resulting weights are positive for
$0<\gamma<\pi$ and $0<v<\pi-\gamma$, with $\Delta=\cos\gamma$.
\item Show that under $u=\epsilon x$, $\eta=\epsilon$,
$\epsilon^{-1}R(u)\to x\id+P$. Identify the three limiting weights.
\end{enumerate}
\end{exercise}

\section{Commuting families}

\begin{exercise}[Monodromy, the RTT relation, and conserved operators]
An $RLL$ relation for operators $L_{0j}(u)$ on an auxiliary copy of $V$
and a site space $W_j$ is
\[
 R_{00'}(u-v)L_{0j}(u)L_{0'j}(v)
 =L_{0'j}(v)L_{0j}(u)R_{00'}(u-v).
\]
Let operators at disjoint sites and auxiliary factors commute, and put
$\mathcal T_0(u)=L_{0L}(u)\cdots L_{01}(u)$ and
$t(u)=\Tr_0\mathcal T_0(u)$. These equations are the defining relations
for the operator families in this exercise.
\begin{enumerate}[label=(\alph*)]
\item Multiply the local relations, keeping the order of factors, to
prove the same relation with $L$ replaced by $\mathcal T$ (the RTT
relation). Use invertibility of $R(u-v)$ and cyclicity in the auxiliary
factors to prove $[t(u),t(v)]=0$, first away from its exceptional
parameters and then by analytic continuation.
\item For the distinguished arrow space define
$A(u),B(u),C(u),D(u)$ by the four input-output kernels of
$\mathcal T_0(u)$ on the auxiliary arrows, so $t(u)=A(u)+D(u)$.
Extract $[B(u),B(v)]=0$ from RTT.
\item Wherever $t(u_0)$ is invertible prove that
$t(u_0)^{-1}t'(u_0)$ commutes with every $t(v)$ and that the
coefficients of the local power series for
$\log(t(u_0)^{-1}t(u))$ commute with one another.
\end{enumerate}
\end{exercise}

\section{Yang--Baxter equation}

\begin{exercise}[The Yang--Baxter identity and the XXZ Hamiltonian]
The Yang--Baxter equation is the operator identity
\[
 R_{12}(u-v)R_{13}(u)R_{23}(v)
 =R_{23}(v)R_{13}(u)R_{12}(u-v).
\]
On the arrow space use $S^z=Z/2$, $S^x=X/2$,
$S^y=Y/2$, where $(Yf)(s)=-\ii s f(-s)$.
\begin{enumerate}[label=(\alph*)]
\item Prove the equation for the trigonometric kernel of the section
Spectral parameter. The sectors of three equal arrows are scalar;
on the sectors with one or two down arrows reduce the identity to
the addition formulas for $\sinh$. Deduce the $RLL$ relation for
$L_{0j}(u)=R_{0j}(u-\theta_j)$ with arbitrary fixed inhomogeneities
$\theta_j$.
\item For $\theta_j=0$ show that $t(0)=(\sinh\eta)^L U$,
where $U$ is cyclic translation in the direction determined by the
ordered monodromy. Differentiate one factor at a time to prove
\[
 t(0)^{-1}t'(0)=
 \frac1{\sinh\eta}\sum_{j=1}^L P_{j,j+1}R'_{j,j+1}(0).
\]
\item Compute the local derivative and deduce that
\[
 H_{\Delta,h}=J\sum_{j=1}^L
 \left(S^x_jS^x_{j+1}+S^y_jS^y_{j+1}
       +\Delta(S^z_jS^z_{j+1}-\tfrac14)\right)
       -h\sum_j(S^z_j-\tfrac12)
\]
commutes with every $t(u)$. State the additive constant and
multiplicative factor relating $H_{\Delta,0}$ to $t(0)^{-1}t'(0)$.
Take $L\ge3$ and periodic indices throughout.
\end{enumerate}
\end{exercise}

\section{Bethe ansatz}

\begin{exercise}[Magnon scattering and periodic Bethe equations]
In $\mathcal H_{L,M}$ a configuration is the set of down-arrow sites
$1\le x_1<\cdots<x_M\le L$. For nonzero complex $z_1,\dots,z_M$
consider the function on configurations
\[
 \psi(x_1,\dots,x_M)=
 \sum_{\pi\in\mathfrak S_M}A_\pi
       \prod_{a=1}^M z_{\pi(a)}^{x_a}.
\]
This is the coordinate Bethe ansatz for the specified lattice
configuration space. Assume all denominators below are nonzero and
the resulting function is not identically zero.
\begin{enumerate}[label=(\alph*)]
\item Away from adjacent down arrows compute the action of
$H_{\Delta,h}$ and obtain
$E=\sum_a[h+\frac J2(z_a+z_a^{-1}-2\Delta)]$.
\item Impose the equation at adjacent sites to derive
\[
 \frac{A_{\pi\circ(a,a+1)}}{A_\pi}
 =S(z_{\pi(a)},z_{\pi(a+1)}),\qquad
 S(z,w)=-\frac{1+zw-2\Delta w}{1+zw-2\Delta z}.
\]
Check $S(z,w)S(w,z)=1$ and compatibility of successive swaps.
Verify directly that the same pairwise conditions handle a string
of three or more adjacent down arrows.
\item Impose periodicity and derive
$z_a^L=\prod_{b\ne a}S(z_b,z_a)$. Prove that a nonzero function
satisfying these conditions is an eigenfunction of $H_{\Delta,h}$.
\end{enumerate}
\end{exercise}

\begin{exercise}[Prediction: the one-magnon band above a saturated chain]
Consider $H_{\Delta,h}$ with $J>0$, $\Delta\ge0$, lattice spacing
$a_0$, and even $L\ge4$. The polarized state has energy zero.
The measured energy variable is $\hbar\omega$ and the wave number is
$k/a_0$, with $k$ defined modulo $2\pi$.
\begin{enumerate}[label=(\alph*)]
\item Use the $M=1$ Bethe equation to derive
\[
 k=\frac{2\pi n}{L},\qquad
 \hbar\omega(k)=h+J(\cos k-\Delta).
\]
Calculate the band width $2J$ and the gap $h-J(1+\Delta)$.
\item Write the exchange on each bond plus half the field on each
endpoint as a positive operator when $h\ge J(1+\Delta)$.
Conclude that the polarized state is a ground state above this
field. For $h<J(1+\Delta)$ use the $k=\pi$ magnon to show that
it is not a ground state.
\item Let $S^-_k=L^{-1/2}\sum_j e^{\ii kj}(S^x_j-\ii S^y_j)$.
Compute its spectral measure in the polarized state: a single
unit-weight line at $\hbar\omega(k)$. Specify the zero-temperature,
nearest-neighbour, and fully polarized assumptions under which
this is a prediction for momentum-resolved spin spectroscopy.
\end{enumerate}
\end{exercise}
